% Clase 2 — Búsqueda Adversaria y Juegos (parte 2)
% Lectura: AIMA Cap. 6
\section{Búsqueda Adversaria y Juegos II}

\subsection{Poda Alfa-Beta ($\alpha$-$\beta$)}

Elimina ramas del árbol que no pueden afectar la decisión final.

\begin{itemize}
    \item $\alpha$: mejor valor que MAX puede garantizar (lo que MAX ya tiene)
    \item $\beta$: mejor valor que MIN puede garantizar (lo que MIN ya tiene)
    \item Poda: si $v \geq \beta$ en nodo MAX $\Rightarrow$ cortar; si $v \leq \alpha$ en nodo MIN $\Rightarrow$ cortar.
\end{itemize}

Con orden perfecto de movimientos: $\bigO{b^{m/2}}$ (duplica la profundidad manejable).

\subsection{Juegos con Incertidumbre: Expectiminimax}

Para juegos con elementos aleatorios (dados) se agregan \textbf{nodos azar}:

$$\text{EXPECTIMINIMAX}(s) = \sum_r P(r) \cdot \text{EXPECTIMINIMAX}(\text{RESULT}(s, r))$$

\subsection{Funciones de Evaluación}

En juegos reales no se puede explorar hasta el final. Se usa una \textbf{función de evaluación heurística} $\text{EVAL}(s)$ que estima la utilidad desde un estado no terminal.

Características de una buena función de evaluación:
\begin{itemize}
    \item Debe concordar con la utilidad real en estados terminales.
    \item Rápida de calcular.
    \item Debe reflejar las probabilidades reales de ganar.
\end{itemize}

\subsection{Monte Carlo Tree Search (MCTS)}

Estrategia basada en simulaciones aleatorias (\textit{rollouts}). Cuatro fases: selección, expansión, simulación, retropropagación.
